% =====================================================================
%  Robot Dynamics (151-0851-00L, HS26) -- MATLAB Cheatsheet
%  ETH-style cheatsheet layout: A4 landscape, multi-column, compact
% =====================================================================
\documentclass[8pt,landscape]{extarticle}

% ---------- Page & fonts ----------
\usepackage[a4paper,landscape,margin=6mm,top=6mm,bottom=7mm]{geometry}
\usepackage[T1]{fontenc}
\usepackage[utf8]{inputenc}
\IfFileExists{lmodern.sty}{\usepackage{lmodern}}{}
\usepackage{amsmath,amssymb,bm}
\usepackage{multicol}
\usepackage{xcolor}
\usepackage{listings}
\usepackage{titlesec}
\usepackage{enumitem}
\usepackage{fancyhdr}
\usepackage[most]{tcolorbox}

% ---------- ETH colours ----------
\definecolor{ethblue}{HTML}{215CAF}
\definecolor{ethpetrol}{HTML}{007894}
\definecolor{ethgreen}{HTML}{627313}
\definecolor{ethred}{HTML}{B7352D}
\definecolor{ethgray}{HTML}{6F6F6F}
\definecolor{codebg}{HTML}{F4F6FA}
\definecolor{mlcomment}{HTML}{228B22}
\definecolor{mlstring}{HTML}{A020F0}
\definecolor{mlkey}{HTML}{0000FF}

% ---------- Layout tweaks ----------
\setlength{\parindent}{0pt}
\setlength{\parskip}{1pt}
\setlength{\columnsep}{3mm}
\setlength{\columnseprule}{0.2pt}
\def\columnseprulecolor{\color{ethgray!50}}
\setlist{nosep,leftmargin=*,label=\textcolor{ethblue}{\scriptsize$\blacktriangleright$}}
\setlength{\abovedisplayskip}{1pt}\setlength{\belowdisplayskip}{1pt}
\setlength{\abovedisplayshortskip}{0pt}\setlength{\belowdisplayshortskip}{0pt}
\renewcommand{\arraystretch}{1.05}

% Section headers: coloured bars (ETH cheatsheet style)
\newcommand{\secbar}[1]{\colorbox{ethblue}{\parbox{\dimexpr\columnwidth-2\fboxsep}{\strut\thesection\ #1}}}
\newcommand{\subbar}[1]{\thesubsection\ #1\ \leaders\hbox{\rule[0.5ex]{1pt}{0.3pt}}\hfill\null}
\titleformat{\section}{\normalfont\bfseries\color{white}}{}{0pt}{\secbar}
\titlespacing*{\section}{0pt}{3pt}{1pt}
\titleformat{\subsection}{\normalfont\bfseries\color{ethblue}}{}{0pt}{\subbar}
\titlespacing*{\subsection}{0pt}{2pt}{0.5pt}
\renewcommand{\thesubsection}{\arabic{section}.\arabic{subsection}}

% Header / footer
\pagestyle{fancy}\fancyhf{}
\renewcommand{\headrulewidth}{0pt}
\fancyfoot[L]{\tiny\color{ethgray}MATLAB Cheatsheet -- Basics, Linear Algebra, Robot Dynamics}
\fancyfoot[C]{\tiny\color{ethgray}Indexing starts at 1 \quad Column vectors by default \quad \texttt{C\_AB}: ${}_A\mathbf r=\mathbf C_{AB}\,{}_B\mathbf r$ \quad quaternion $[w;x;y;z]$}
\fancyfoot[R]{\tiny\color{ethgray}Aaryan -- page \thepage}

% ---------- MATLAB listings ----------
\lstdefinestyle{ml}{
  language=Matlab,
  basicstyle=\ttfamily\fontsize{6.2}{7}\selectfont,
  keywordstyle=\color{mlkey},
  commentstyle=\color{mlcomment}\itshape,
  stringstyle=\color{mlstring},
  backgroundcolor=\color{codebg},
  frame=none, 
  xleftmargin=3pt, framexleftmargin=2pt,
  aboveskip=1pt, belowskip=1pt,
  columns=fullflexible, keepspaces=true,
  breaklines=true, breakatwhitespace=false,
  showstringspaces=false, upquote=true,
  morekeywords={syms,sym,jacobian,simplify,matlabFunction,subs,diff,pinv,null,ode45,quadprog,assert,feval,skew,sgn},
  literate={~}{{$\sim$}}1
}
\lstnewenvironment{ml}{\lstset{style=ml}}{}
\newcommand{\mc}[1]{\texttt{\detokenize{#1}}}

% Warning/Tip boxes
\newtcolorbox{tip}{colback=ethgreen!8,colframe=ethgreen,boxrule=0.4pt,left=1pt,right=1pt,top=0pt,bottom=0pt,arc=0pt,fontupper=\scriptsize}
\newtcolorbox{warn}{colback=ethred!7,colframe=ethred,boxrule=0.4pt,left=1pt,right=1pt,top=0pt,bottom=0pt,arc=0pt,fontupper=\scriptsize}

% Math shortcuts
\newcommand{\C}{\mathbf{C}}
\newcommand{\T}{\mathbf{T}}
\newcommand{\J}{\mathbf{J}}
\newcommand{\M}{\mathbf{M}}
\newcommand{\q}{\mathbf{q}}
\newcommand{\dq}{\dot{\mathbf{q}}}
\newcommand{\ddq}{\ddot{\mathbf{q}}}
\newcommand{\bb}{\mathbf{b}}
\newcommand{\g}{\mathbf{g}}
\newcommand{\bI}{\mathbb{I}}
\newcommand{\skw}[1]{[#1]_\times}
\newcommand{\pinv}{^{+}}

\begin{document}
\scriptsize
\begin{multicols*}{4}

\begin{tcolorbox}[colback=ethblue,colframe=ethblue,coltext=white,arc=0pt,left=2pt,right=2pt,top=1pt,bottom=1pt]
{\bfseries\normalsize MATLAB Cheatsheet}\\
{\tiny Basics \textbullet{} Linear Algebra \textbullet{} Robot Dynamics (151-0851-00L, HS26)}
\end{tcolorbox}

% =====================================================================
\section{Environment \& Workspace}
\begin{ml}
clc                 % clear command window
clear / clear x     % clear all / one variable
close all           % close all figures
clearvars -except q % keep q
who, whos           % list variables (with size/class)
help pinv           % short help;  doc pinv -> full docs
which myFun         % which file is called?
exist('x','var')    % 1 if variable x exists
pwd, cd('..'), ls   % folders
addpath(genpath('ex1'))   % add folder + subfolders to path
format long / format short / format compact
x = 3;              % ; suppresses output
x = 3               % no ; -> prints x
...                 % line continuation
% comment    %% section (Ctrl+Enter runs section)
Ctrl+C              % abort running code
\end{ml}
\begin{tip}Constants: \mc{pi}, \mc{Inf}, \mc{NaN}, \mc{eps} ($2.2\cdot10^{-16}$), \mc{1i} (imag.\ unit). \mc{ans} = last unassigned result.\end{tip}

% =====================================================================
\section{Creating Arrays}
\begin{ml}
a = 5;                    % scalar (1x1 double)
r = [1 2 3];  r = [1, 2, 3];   % row vector 1x3
c = [1; 2; 3];            % column vector 3x1
c = r';                   % transpose -> column
A = [1 2; 3 4];           % 2x2 matrix (; = new row)
x = 0:0.1:1;              % start:step:end
x = 1:6;                  % step 1
x = 10:-2:0;              % descending
x = linspace(0, 1, 50);   % 50 points incl. ends
zeros(3), zeros(3,6), ones(2,3), eye(4)
rand(3,1)   % U(0,1)    randn(6,1)  % N(0,1)
randi(10,1,5)             % 5 ints in 1..10
diag([1 2 3])             % diagonal matrix
diag(A)                   % extract diagonal (column)
blkdiag(A, B)             % block-diagonal
repmat(v, 1, 4)           % tile
[A, B]  / horzcat(A,B)    % side by side
[A; B]  / vertcat(A,B)    % stacked
cat(3, A, B)              % 3-D array
[X,Y] = meshgrid(x, y);   % grids for surf/contour
NaN(3,1), true(2), false(1,4)
\end{ml}

% =====================================================================
\section{Size \& Indexing}
\begin{ml}
size(A)        % [rows cols]      size(A,1) rows
numel(A)       % # elements       length(v) longest dim
ndims(A), isempty(A), isrow(v), iscolumn(v)
A(2,3)         % row 2, col 3
A(2,:)         % whole row 2      A(:,end) last column
A(1:2, [1 3])  % sub-block
A(end, end)    % last element
A(:)           % all elements as column (column-major!)
A(5)           % linear index (column-wise counting)
A(:,2) = [];   % delete column 2
v(end+1) = 7;  % append
A(A < 0) = 0;  % logical indexing
idx = find(v > 2);          % indices where true
[r, c] = find(A == max(A(:)));
k = sub2ind(size(A), r, c); [r,c] = ind2sub(size(A), k);
B = reshape(A, 3, []);      % [] = auto size
v = v(:);                   % force column
flip(v), fliplr(A), flipud(A), circshift(v, 1)
\end{ml}
\begin{warn}MATLAB stores \textbf{column-major}: \mc{A(:)} stacks columns. Indices start at \textbf{1}. \mc{T(1:3,1:3)} = rotation, \mc{T(1:3,4)} = translation.\end{warn}

% =====================================================================
\section{Operators}
\begin{ml}
+  -             % add/sub (implicit expansion: A - v')
*  /  \  ^       % MATRIX mult, right/left divide, power
.* ./ .^         % ELEMENT-WISE
'                % conjugate transpose (use .' for sym/cplx)
== ~= < <= > >=  % relational -> logical array
&  |  ~  xor     % element-wise logical
&& ||            % short-circuit (scalars, in if/while)
any(v), all(v)   % any/all true
isequal(A, B)    % exact equality (shape + values)
abs(a-b) < 1e-9  % float comparison with tolerance
\end{ml}
\begin{tip}\mc{A*B} needs inner dims to match; \mc{A.*B} needs equal (or broadcastable) sizes. \mc{A^2 = A*A} but \mc{A.^2} squares entries.\end{tip}

% =====================================================================
\section{Elementary Functions}
\begin{ml}
sin cos tan asin acos atan  atan2(y,x)  % radians!
sind cosd  deg2rad rad2deg
exp log log10 sqrt abs sign mod rem
round(x,3)  floor ceil fix
sum(A)       % column sums      sum(A,2)  row sums
sum(A,'all') prod cumsum cumprod diff(v)
mean median std var  min max
[m, i] = max(v);          % value and index
max(A, [], 'all'), max(A, 0)   % elementwise max with 0
sort(v,'descend'); [s, i] = sort(v);
unique(v), ismember(3, v), numel(find(v>0))
norm(v)  vecnorm(A)       % column-wise 2-norms
\end{ml}

% =====================================================================
\section{Control Flow}
\begin{ml}
if x > 0 && ~isempty(v)
    ...
elseif x == 0
    ...
else
    ...
end
for k = 1:n            % also: for col = A (iterates cols)
    if k == 3, continue; end
    if done,   break;    end
end
while norm(e) > tol && it < maxIt
    it = it + 1;
end
switch mode
    case 'pos',        ...
    case {'rot','ori'},...
    otherwise,         ...
end
try
    x = riskyFun();
catch err
    disp(err.message);
end
\end{ml}

% =====================================================================
\section{Functions}
\subsection{Function file \texttt{myFun.m} (name = file name)}
\begin{ml}
function [y, J] = myFun(q, k)
%MYFUN  One-line description (shown by help)
    if nargin < 2, k = 1; end      % default argument
    y = k*q;
    if nargout > 1, J = k*eye(numel(q)); end
end

function h = helper(x)   % local function: same file,
    h = 2*x;             % after the main one, private
end
\end{ml}
Scripts: plain \mc{.m} without \mc{function}, share the base workspace.
\subsection{Anonymous functions \& handles}
\begin{ml}
f  = @(x) x.^2 + 3*x;          % f(2) -> 10
g  = @(x,y) x'*y;
Rz = @(a) [cos(a) -sin(a) 0; sin(a) cos(a) 0; 0 0 1];
h  = @myFun;   h(q)            % handle to file fn
feval('myFun', q)              % call by name string
fplot(f, [-2 2])
% anonymous fns capture variables at creation time!
\end{ml}
\subsection{Variable arguments \& map-style calls}
\begin{ml}
function out = f(varargin)      % varargin{1}, ...
    n = numel(varargin);
end
[a, ~] = myFun(q);              % ignore output
arrayfun(@(k) k^2, 1:5)         % apply to each element
arrayfun(@(k) Rz(k), [0 pi], 'UniformOutput', false)
cellfun(@numel, {'ab','cde'})   % -> [2 3]
\end{ml}

% =====================================================================
\section{Data Types}
\begin{ml}
class(x), isa(x,'double'), isnumeric, islogical
double(x), int32(x), logical(x)
% --- char vs string
s1 = 'text';  s2 = "text";      % char array vs string
[s1 ' more'],  s2 + " more",  strcat('a','b')
num2str(pi,4), str2double('3.1'), mat2str(A)
sprintf('q%d', 3)               % -> 'q3'
strcmp(a, b), contains(s, 'x'), upper(s), split(s, ',')
% --- cell arrays (heterogeneous containers)
c = {1, 'two', [3 4 5]};
c{3}          % content -> [3 4 5]
c(3)          % sub-cell -> {[3 4 5]}
c = cell(1,6);  c{k} = T_k;     % e.g. store transforms
% --- structs
p.m = 2.5;  p.l = [0.3 0.2];  p.name = 'link1';
p = struct('m', 2.5, 'l', 0.3);
p.('m'), fieldnames(p), isfield(p,'m'), rmfield(p,'m')
links(2).m = 1.2;               % struct array
[links.m]                       % collect field values
% --- maps
M = containers.Map({'a','b'}, {1, 2});  M('a')
\end{ml}

% =====================================================================
\section{Printing \& Files}
\begin{ml}
disp(A)                    disp(['q = ', num2str(q')])
fprintf('err = %.3e, it = %d\n', err, it);
fprintf('%6.3f %6.3f %6.3f\n', A');   % prints ROWS of A
% %d int, %f float, %.4f, %e sci, %g compact, %s string
save('data.mat', 'q', 'T');  load('data.mat');
S = load('data.mat');  S.q
writematrix(A, 'A.csv');  A = readmatrix('A.csv');
T = readtable('log.csv');  T.time
\end{ml}

% =====================================================================
\section{Linear Algebra}
\subsection{Products \& basic quantities}
\begin{ml}
A*B          % matrix product     a'*b   % dot (cols)
dot(a,b)     cross(a,b)           % 3-vectors
a*b'         % outer product      kron(A,B)
A'  A.'      % transposes         A^k    % matrix power
trace(A), det(A), rank(A), rank(A, tol)
norm(v)      % 2-norm     norm(v,1), norm(v,Inf)
norm(A)      % spectral   norm(A,'fro')
cond(A)      % condition number (large -> near singular)
triu(A), tril(A), diag(A), diag(diag(A))
issymmetric(A), isdiag(A), istriu(A)
\end{ml}
\subsection{Solving linear systems}
\begin{ml}
x = A \ b;        % solve A x = b  (DON'T use inv(A)*b)
X = B / A;        % solve X A = B
x = pinv(A) * b;  % Moore-Penrose: least-squares + MIN-NORM
x = lsqminnorm(A, b);   % min-norm LS, faster than pinv
Ainv = inv(A);    % only if you really need the matrix
\end{ml}
\begin{warn}\texttt{A\textbackslash b} for a \textbf{fat} (underdetermined) $A$ returns a \emph{basic} solution with many zeros -- \textbf{not} the minimum-norm one. For redundant robots ($\mathbf J\in\mathbb R^{m\times n}$, $n>m$) use \mc{pinv(J)} / \mc{lsqminnorm}. For tall $A$, \texttt{A\textbackslash b} = least squares.\end{warn}
\subsection{Decompositions}
\begin{ml}
e       = eig(A);           % eigenvalues
[V, D]  = eig(A);           % A*V = V*D
s       = svd(A);           % singular values (descending)
[U,S,V] = svd(A);           % A = U*S*V'
[U,S,V] = svd(A, 'econ');   % economy size
[Q, R]  = qr(A);            % A = Q*R
[L,U,P] = lu(A);            % P*A = L*U
R       = chol(A);          % A = R'*R  (A must be SPD)
[R, p]  = chol(A);          % p == 0  <=>  A is SPD
N       = null(A);          % orthonormal basis of ker(A)
O       = orth(A);          % orthonormal basis of range(A)
expm(A), logm(A), sqrtm(A)  % matrix exp/log/sqrt
\end{ml}
\subsection{Useful checks}
\begin{ml}
all(eig((M+M')/2) > 0)          % positive definite?
norm(C'*C - eye(3)) < 1e-9      % orthogonal?
abs(det(C) - 1) < 1e-9          % proper rotation?
min(svd(J))                     % distance to singularity
sqrt(det(J*J'))                 % manipulability
\end{ml}

% =====================================================================
\section{Symbolic Math Toolbox}
\begin{ml}
syms x y real                     % scalar symbols
syms q [6 1] real                 % q1..q6 as vector q
q = sym('q', [6 1], 'real');      % same, older versions
syms m l g positive
assume(x > 0)
f   = x^2*sin(y);
df  = diff(f, x);   d2f = diff(f, x, 2);
J   = jacobian([f; x*y], [x; y]); % Jacobian of vector fn
H   = hessian(f, [x; y]);
F   = int(f, x);    F = int(f, x, 0, 1);
f2  = simplify(f);  expand(f), collect(f, x), factor(f)
subs(f, x, 2)                     % substitute
subs(f, {x, y}, {1, pi})
double(subs(f, [x y], [1 2]))     % -> numeric
vpa(pi, 8), pretty(f), latex(f)
sol = solve(x^2 == 4, x);         % symbolic solve
[sx, sy] = solve([x+y==3, x-y==1], [x y]);
taylor(sin(x), x, 'Order', 4)
A = sym(zeros(3));                % symbolic preallocation
% symbolic -> fast numeric function / m-file
fh = matlabFunction(J, 'Vars', {x, y});
matlabFunction(M, 'File','M_fun', 'Vars', {q}, ...
               'Optimize', false);  % faster generation
\end{ml}
\begin{tip}Time derivative of a matrix $\mathbf A(\mathbf q)$: \mc{dA = reshape(jacobian(A(:),q)*dq, size(A))}. Avoid \mc{simplify} on huge expressions -- it can take minutes. Use \mc{.'} (not \mc{'}) or declare symbols \mc{real}, else \mc{'} conjugates.\end{tip}

% =====================================================================
\section{Numerical Methods}
\begin{ml}
% ODE:  dx/dt = f(t,x)
opts = odeset('RelTol',1e-8, 'AbsTol',1e-10);
[t, X] = ode45(@(t,x) f(t,x,p), [0 5], x0, opts);
% X: rows = time steps, cols = states;  X(:,1) = x1(t)
[t, X] = ode45(fh, 0:0.01:5, x0);     % fixed output grid
% stiff: ode15s;  fixed step: write your own loop
for k = 1:N, x = x + dt*f(t(k), x); end   % expl. Euler

xr = fzero(@(x) cos(x) - x, 0.5);     % scalar root
xs = fsolve(@(x) F(x), x0);           % nonlinear system
xm = fminsearch(@(x) cost(x), x0);    % derivative-free
xm = fminunc(@cost, x0);              % Optim. TB
x  = quadprog(H, f, A, b, Aeq, beq, lb, ub);  % QP
x  = lsqlin(C, d, A, b, Aeq, beq);    % constrained LS
yi = interp1(t, y, ti, 'linear');     % interpolation
I  = integral(@(x) x.^2, 0, 1);
trapz(t, y), cumtrapz(t, y), gradient(y, dt)
pp = polyfit(x, y, 3);  polyval(pp, xq);  roots(pp)
\end{ml}
\subsection{Control System Toolbox}
\begin{ml}
sys = ss(A, B, C, D);  tf(sys),  step(sys), bode(sys)
K = lqr(A, B, Q, R);   eig(A - B*K)
K = place(A, B, poles);
rank(ctrb(A, B)), rank(obsv(A, C))
\end{ml}

% =====================================================================
\section{Plotting}
\begin{ml}
figure('Name','Joints');  hold on;  grid on;
plot(t, q(1,:), 'b-', 'LineWidth', 1.5);
plot(t, q_d(1,:), 'r--');
xlabel('t [s]'); ylabel('q_1 [rad]'); title('Tracking');
legend('q','q_d','Location','best');  xlim([0 5]);
subplot(2,1,1); ...  subplot(2,1,2); ...
tiledlayout(2,2); nexttile; plot(...);
plot3(x, y, z); axis equal; view(3);  % 3-D path
scatter(x, y, 10, c, 'filled');
quiver3(x0,y0,z0, u,v,w, 0);           % frame axes/arrows
surf(X, Y, Z); shading interp; colorbar;
yline(0,'k:'); xline(1);
exportgraphics(gcf, 'fig.pdf');        % save figure
% animation
h = plot3(nan,nan,nan,'o-');
for k = 1:numel(t)
    set(h, 'XData', P(1,:,k), 'YData', P(2,:,k), 'ZData', P(3,:,k));
    drawnow;  % pause(0.01)
end
\end{ml}

% =====================================================================
\section{Robot Dynamics Toolkit}
\subsection{Rotations \& transforms}
\begin{ml}
skew = @(v) [0 -v(3) v(2); v(3) 0 -v(1); -v(2) v(1) 0];
unskew = @(S) [S(3,2); S(1,3); S(2,1)];
Rx = @(a) [1 0 0; 0 cos(a) -sin(a); 0 sin(a) cos(a)];
Ry = @(a) [cos(a) 0 sin(a); 0 1 0; -sin(a) 0 cos(a)];
Rz = @(a) [cos(a) -sin(a) 0; sin(a) cos(a) 0; 0 0 1];
C_IB  = Rz(yaw)*Ry(pitch)*Rx(roll);   % ZYX Euler
A_r   = C_AB * B_r;                   % change frame
C_AC  = C_AB * C_BC;                  % chain
C_BA  = C_AB';                        % inverse
% homogeneous transforms
T  = [C, r; 0 0 0 1];
Ti = [C', -C'*r; 0 0 0 1];            % inverse
p  = T * [r_B; 1];  p = p(1:3);        % map point
C  = T(1:3,1:3);   r = T(1:3,4);
% rotation vector <-> matrix via matrix exp/log
C   = expm(skew(phi));                % phi = theta*n
phi = unskew(logm(C));                % (theta < pi)
th  = acos((trace(C)-1)/2);           % rotation angle
% Euler angles from C (ZYX): use atan2, never atan
yaw   = atan2(C(2,1), C(1,1));
pitch = atan2(-C(3,1), sqrt(C(3,2)^2 + C(3,3)^2));
roll  = atan2(C(3,2), C(3,3));
wrap  = @(a) atan2(sin(a), cos(a));   % to (-pi,pi]
\end{ml}
\subsection{Quaternions $[w;\,x;\,y;\,z]$}
\begin{ml}
qmul = @(p,q) [p(1)*q(1) - p(2:4)'*q(2:4);
               p(1)*q(2:4) + q(1)*p(2:4) + cross(p(2:4),q(2:4))];
qconj = @(q) [q(1); -q(2:4)];
q = q / norm(q);                      % normalise!
q2C = @(q) (2*q(1)^2-1)*eye(3) + 2*q(1)*skew(q(2:4)) ...
           + 2*q(2:4)*q(2:4)';
\end{ml}
\begin{tip}Robotics/Navigation TB built-ins: \mc{eul2rotm([z y x])} (default \mc{'ZYX'}), \mc{rotm2eul}, \mc{rotm2quat} (returns \emph{row} \mc{[w x y z]}), \mc{quat2rotm}, \mc{axang2rotm([n' th])}. Cross-check against your own functions -- the exercises auto-grade \emph{your} implementation.\end{tip}
\subsection{Loops over a kinematic chain}
\begin{ml}
T_I = cell(1,n);  T = T_I0;
for k = 1:n
    T = T * T_k(q, k);            % T_{k-1,k}(q_k)
    T_I{k} = T;                   % store T_Ik
end
r_IE = T_I{n}(1:3,4);
% geometric Jacobian, revolute joints (all in I)
J_P = zeros(3,n); J_R = zeros(3,n);
for k = 1:n
    n_k = T_I{k}(1:3,1:3) * axis_k(:,k);    % axis in I
    J_P(:,k) = cross(n_k, r_IE - T_I{k}(1:3,4));
    J_R(:,k) = n_k;
end
J = [J_P; J_R];
\end{ml}
\subsection{Finite-difference check of any Jacobian}
\begin{ml}
function J = numJac(f, q, h)
    if nargin < 3, h = 1e-6; end
    f0 = f(q); J = zeros(numel(f0), numel(q));
    for k = 1:numel(q)
        e = zeros(size(q)); e(k) = h;
        J(:,k) = (f(q+e) - f(q-e)) / (2*h);
    end
end
norm(numJac(@jointToPosition, q) - jointToPosJac(q))
\end{ml}
\subsection{Pseudo-inverse \& null-space}
\begin{ml}
Jp   = pinv(J);                               % J^+
Jd   = J' / (J*J' + lambda^2*eye(size(J,1))); % damped
N    = eye(n) - pinv(J)*J;                    % projector
dq   = pinv(J1)*w1 + N1*pinv(J2*N1)*(w2 - J2*pinv(J1)*w1);
\end{ml}
\subsection{Equations of motion}
\begin{ml}
% M(q) ddq + b(q,dq) + g(q) = tau
ddq = M \ (tau - b - g);                   % forward dyn.
tau = M*ddq_des + b + g;                   % inverse dyn.
Lam = inv(J/M*J');                         % op-space inertia
% simulate
rhs = @(t,x) [x(n+1:end);
      M_fun(x(1:n)) \ (tau_fun(t,x) - b_fun(x(1:n),x(n+1:end)) - g_fun(x(1:n)))];
[t, X] = ode45(rhs, [0 5], [q0; dq0]);
\end{ml}

% =====================================================================
\section{Debugging \& Performance}
\begin{ml}
dbstop if error      % break into debugger on error
keyboard             % manual breakpoint in code
dbcont / dbquit      % continue / leave debugger
size(x), class(x)    % first thing to check on errors
assert(isequal(size(J), [6 6]), 'J has wrong size');
tic; ...; toc        % timing
profile on; run; profile viewer
q = zeros(6,N);      % PREALLOCATE before loops
y = A*x;             % vectorise instead of loops
\end{ml}
\begin{warn}\textbf{Common errors:} ``Dimensions of arrays being concatenated are not consistent'' $\to$ row/column mix. ``Incorrect dimensions for matrix multiplication'' $\to$ use \mc{.*} or transpose. ``Index exceeds array bounds'' $\to$ check \mc{size}. ``Undefined function'' $\to$ path/\mc{addpath}, file name $\neq$ function name.\end{warn}

\end{multicols*}
\end{document}
